十多年前，我的人工智能之路始于统计模式分类。那时我学到，只要模型能捕捉到数据中有用的结构，相对简单的模型也能走得很远。随着时间推移，我的兴趣从识别模式的模型，转向了能够解释、规划和求解多步问题的模型。

2022 年 ChatGPT 发布后，大型语言模型（LLM）走进了大众视野。随后的一次转变——在 2024 年底到 2025 年期间尤其明显——是 LLM 的延伸，即创造所谓的推理模型，它们能够求解更复杂的问题，尤其是在数学和代码等技术领域。

与此同时，推理模型又常常显得不透明，而本书正是我尝试让推理模型变得更容易上手的一次努力。我们不会从零训练一个庞大的模型，而是从一个预训练的小型 LLM 出发，一步步应用推理技术。在此过程中，你将实现文本生成、用验证器进行评估、推理时扩展方法，以及训练方法，例如可验证奖励的强化学习和从零开始的蒸馏。读到最后，你将理解各种主要推理方法在实践中如何运作，以及它们如何融入现代的 LLM 开发工作流。

与我此前出版的\emph{《从零构建大语言模型》}一样，本书也采用代码优先的方式，但侧重点不同。本书不深入讲解 Transformer 架构，也不涉及大规模预训练，而是聚焦于那些把常规 LLM 变成推理模型的方法。如果你读过前一本书，这一本读起来应当像是自然的延续。如果没有读过，你依然可以跟得上，不必事先了解全部底层 LLM 细节。

全书都会使用数学示例，因为这类示例在推理研究中既实用，又便于自动验证。这使它们成为理解模型是否真正正确求解问题的有用试验场。更宽泛地说，我的目标不是教你应对某一个狭窄的基准测试，而是给你在自己的工作中构建、评估和改进推理系统所需的工具与直觉。

我刚开始研究推理模型时，不得不从论文、代码仓库和零散的实现笔记中拼凑各种想法。我希望本书能为你省下这番工夫，给出一条从第一性原理通向可运行代码的更清晰路径。我坚信，理解推理模型最好的方式就是亲手构建一个。

祝阅读愉快，编码愉快！
